Fact \claimref{odd-extensions.imported-frozen-point} gives the signature of $I^{\perp}$ as $(1,4)$ or $(2,3)$ by the sign of $\Lambda$, the sign of $Q(V) = -3/\Lambda$, and $\{I,V\}^{\perp}$ of signature $(1,3)$ in both cases; so $e_{5}^{2} = \operatorname{sign}Q(V)$ is $-1$ for $\Lambda > 0$ and $+1$ for $\Lambda < 0$, and $\Cl(I^{\perp})$ is the space extension or the time extension of $\Cl(1,3)$ respectively. The translation part of the reduced connection is $-\hb^{-1}(a)\wedge V/\ell^{2} = -\hb^{-1}(a)\wedge e_{5}/\ell$ ($Q(V) = -\ell^{2}$, $V = \ell e_{5}$), whose Clifford image is the bivector $e_{\mu}e_{5}$ up to the scalar $-1/\ell$; the algebra these generate is Theorem \claimref{odd-extensions.mirror-inside-complexified-extension-absent-doubled}: $\Cl(1,3)$ in the first case, $\Cl(3,1)$ in the second. The contraction rescales the generators by $1/\ell$ and rescaling does not change the signs of the squares.
